%杨舒云的实验报告编辑界面，使用了Huanyu Shi，2019级的模板，杨舒云在此拜谢ORZ!

%!TEX program = xelatex
\documentclass[dvipsnames, svgnames,a4paper,11pt]{article}
\input{YsySettings} 
\usepackage{lipsum}
\usepackage{adjustbox}

\begin{document}
	
	
	
	%实验报告封面	
	\begin{table}
		\renewcommand\arraystretch{1.7}
		\begin{tabularx}{\textwidth}{
				|X|X|X|X
				|X|X|X|X|}
			\hline
			\multicolumn{2}{|c|}{预习报告}&\multicolumn{2}{|c|}{实验记录}&\multicolumn{2}{|c|}{分析讨论}&\multicolumn{2}{|c|}{总成绩}\\
			\hline
			\LARGE20 & & \LARGE30 & & \LARGE30 & & \LARGE80 & \\
			\hline
		\end{tabularx}
	\end{table}
	
	\begin{table}
		\renewcommand\arraystretch{1.7}
		\begin{tabularx}{\textwidth}{|X|X|X|X|}
			\hline
			年级、专业： & 2022级 物理学 &组号： & 2\\
			\hline
			姓名： & 杨舒云  & 学号： & 22344020\\
			\hline
			实验时间： & 2023/12/7 & 教师签名： & \\
			\hline
		\end{tabularx}
	\end{table}
	
	\begin{center}
		\LARGE AC4 \quad 冰的熔化热测量
	\end{center}
	
	\textbf{【实验报告注意事项】}
	\begin{enumerate}
		\item 实验报告由三部分组成：
		\begin{enumerate}
			\item 预习报告：课前认真研读实验讲义，弄清实验原理；实验所需的仪器设备、用具及其使用、完成课前预习思考题；了解实验需要测量的物理量，并根据要求提前准备实验记录表格（可以参考实验报告模板，可以打印）。\textcolor{red}{\textbf{（25分）}}
			\item 实验记录：认真、客观记录实验条件、实验过程中的现象以及数据。实验记录请用珠笔或者钢笔书写并签名（\textcolor{red}{\textbf{用铅笔记录的被认为无效}}）。\textcolor{red}{\textbf{保持原始记录，包括写错删除部分，如因误记需要修改记录，必须按规范修改。}}（不得输入电脑打印，但可扫描手记后打印扫描件）；离开前请实验教师检查记录并签名。\textcolor{red}{\textbf{（30分）}}
			\item 数据处理及分析讨论：处理实验原始数据（学习仪器使用类型的实验除外），对数据的可靠性和合理性进行分析；按规范呈现数据和结果（图、表），包括数据、图表按顺序编号及其引用；分析物理现象（含回答实验思考题，写出问题思考过程，必要时按规范引用数据）；最后得出结论。\textcolor{red}{\textbf{（25分）}}
		\end{enumerate}
		\textbf{实验报告就是将预习报告、实验记录、和数据处理与分析合起来，加上本页封面。\textcolor{red}{（80分）}}
		\item 每次完成实验后的一周内交\textbf{实验报告}（特殊情况不能超过两周）。
	\end{enumerate}
	
	\textbf{【实验安全注意事项】}	
	\begin{enumerate}
		\item 使用热水和冰块注意避免烫伤、冻伤。
		\item 运输水时，注意他人的安全，避免把水弄到地面上。	
	\end{enumerate}
	
	\textbf{【特别鸣谢及模板说明】}	
	
	感谢2019级学长石寰宇为本实验报告提供\LaTeX 模板。
	
	
	
	\clearpage
	\tableofcontents
	\clearpage
	
	
	
	%预习报告	
	\setcounter{section}{0}
	\section{AC4 冰的熔化热测量 \quad\heiti 预习报告}
	
	\subsection{实验目的}
	\begin{enumerate}
		\item 掌握混合法测量冰的熔化热基本原理，学习物理建模；
		\item 测定冰的熔化热。
	\end{enumerate}
	
	\subsection{仪器用具}
	\begin{table}[htbp]
		\centering
		\renewcommand\arraystretch{1.6}
		% \setlength{\tabcolsep}{10mm}
		\begin{tabular}{p{0.05\textwidth}|p{0.20\textwidth}|p{0.05\textwidth}|p{0.5\textwidth}}
			\hline
			编号& 仪器用具名称 & 数量 &  主要参数（型号，测量范围，测量精度等） \\
			\hline
			1& MS6514测温仪 & 2 & 配两对K型热电偶 \\
			
			2& 保温杯 & 1 & 350ml，或500ml \\
			
			3& 量杯 & 1 & 量程1000g，最小分辨率0.01g \\
			
			4& 电子天平 & 1 & -- \\
			\hline
		\end{tabular}
	\end{table}
	
	
	
	\subsection{原理概述}
	引言：温度测量和量热技术是热学实验中最基本的问题。量热学是以热力学第一定律（能量守恒定律） 为理论基础的，它所研究的范围就是如何计量物质系统随温度变化、相变、化学反应等吸收或放出的热量。量热学的常用实验方法有混合法、稳流法、冷却法、替热法、电热法等。本实验学习利用混合法来测定冰的熔化热； 使用的基本仪器为量热器。 本实验用绝热效果好的家用保温杯取代量热器， 虽然减少了热量散失对实验的影响， 但不能直接测量与液体接触的保温杯内胆壁的质量， 这使本实验为学生测量该热质量提供了研究空间。
	
	\subsubsection{简述热力学第一定律}
	
	热力学第一定律是能量守恒和转换的基本原理，应用于热力学系统。它可以表述为：系统内能的变化等于系统所吸收的热量与系统对外做的功之和。
	
	数学上，这可以表示为：
	\[ \Delta U = Q - W \]
	
	其中：\(\Delta U\)是系统内能的变化量；\(Q\)是系统吸收的热量；\(W\)是系统对外做的功。
	
	这个公式有几个要点需要注意：内能（U）：是系统内部所有微观粒子（如分子、原子）动能和势能的总和，包括化学能、核能等。热量（Q）：是能量的一种传递形式，当两个温度不同的物体接触时，能量会从温度高的物体流向温度低的物体。功（W）：在热力学中，通常指由系统对外界做的机械功，例如气体膨胀对外界做功。
	
	热力学第一定律的核心含义是能量守恒。在一个孤立系统中，能量不能被创造或消灭，只能从一种形式转换为另一种形式，或者从一个物体传递到另一个物体。这意味着任何能量的增加或减少都必须有一个相应的能量输入或输出。

	\begin{figure}[htbp]
		\centering
		\includegraphics[width=0.5\textwidth]{AC4GraAD1.png}
		\caption{热力学第一定律}
		\label{fig:figAD1}
	\end{figure}
	
	\subsubsection{简述利用混合法来测定冰的熔化热的原理，列出必要的公式}
	利用混合法测定冰的熔化热是一种经典的热力学实验方法。其基本原理是基于热量守恒，即在一个绝热系统中，热量的损失和获得总是平衡的。
	\begin{figure}[htbp]
		\centering
		\includegraphics[width=0.5\textwidth]{AC4GraAD2.png}
		\caption{利用混合法来测定冰的熔化热的原理示意图}
		\label{fig:figAD2}
	\end{figure}
	
	利用混合法测定冰的熔化热是基于热量守恒的原理，通常包括以下几个步骤：
	\begin{enumerate}
		\item 系统设置：通常包括一个绝热容器（如卡尔文杯），内含已知质量和温度的水。将质量和温度已知的冰块加入到水中。
		\item 热量交换：冰块吸收热量融化成水，而容器中的水则失去相同的热量。整个系统是绝热的，意味着没有热量与外界交换。
		\item 温度变化：混合后，系统达到热平衡，此时系统的温度介于初始水温和冰的温度之间。
		\item 热量守恒：在这个过程中，水失去的热量等于冰吸收的热量加上冰融化时吸收的熔化热。
	\end{enumerate}
	
	为了计算冰的熔化热（\( L \)），我们可以使用以下公式：
	
	\begin{itemize}
		\item 水失去的热量:
		\[ Q_{\text{水失去}} = m_{\text{水}} c_{\text{水}} (T_{\text{初水}} - T_{\text{终}}) \]
		其中 \( m_{\text{水}} \) 是水的质量，\( c_{\text{水}} \) 是水的比热容，\( T_{\text{初水}} \) 是初始水温，\( T_{\text{终}} \) 是最终温度。
		
		\item 冰吸收的热量（未包括熔化热）:
		\[ Q_{\text{冰吸收}} = m_{\text{冰}} c_{\text{冰}} (T_{\text{终}} - T_{\text{初冰}}) \]
		其中 \( m_{\text{冰}} \) 是冰的质量，\( c_{\text{冰}} \) 是冰的比热容，\( T_{\text{初冰}} \) 是初始冰温（通常是 \( 0^\circ C \)）。
		
		\item 冰的熔化热:
		\[ Q_{\text{熔化}} = m_{\text{冰}} L \]
		其中 \( L \) 是冰的熔化热。
		
		\item 热量守恒:
		\[ Q_{\text{水失去}} = Q_{\text{冰吸收}} + Q_{\text{熔化}} \]
	\end{itemize}
	
	结合以上公式，可以解出冰的熔化热 \( L \)：
	\[ L = \frac{m_{\text{水}} c_{\text{水}} (T_{\text{初水}} - T_{\text{终}}) - m_{\text{冰}} c_{\text{冰}} (T_{\text{终}} - T_{\text{初冰}})}{m_{\text{冰}}} \]
	
	
	\subsection{实验前思考题}
	\begin{question}
		讲义中公式$mL+mC_0(T_1-T')=(m_0C_0+m_1C_1)(T_0-T_1)$(AC4- 1)的建模做了哪些简化？
	\end{question}
	\begin{enumerate}
		\item 忽略系统与外界的热交换：在推导过程中，假设量热器（即实验系统）与外界没有热交换，将水、冰和量热器看作是孤立系统。这意味着所有的热量交换只发生在系统的内部，即冰熔化成水并升温所吸收的热量等于水和量热器内筒降温所放出的热量。
		
		\item 假定冰块完全熔化：在计算中，假设加入的冰块在达到热平衡时已完全熔化成水。这意味着冰块在过程中完全转化为水，且在这个过程中吸收了相应的熔化热。
		
		\item 忽略温度变化带来的其他影响：在实际情况中，由于热交换的存在，系统在投入冰块前的温度以及冰块完全熔化后的温度并非定值。这个简化假设在某些情况下可能导致实验精度的降低，因此实验设计中需要减少漏热或修正漏热的影响，以保证测量方案的准确性。
		
		\item 除此之外，公式中没有考虑搅拌器的吸热，即假设搅拌器的质量与比热远小于水和量热器，其吸的热相比之下可忽略。假设系统是绝热的，热量不会损失到环境。假设冰和水在混合后迅速达到热平衡态，忽略达到平衡所需的时间。
	\end{enumerate}
	
	\begin{question}
		保温杯$m_1C_1$估算所基于的模型作了什么简化？请分析其误差。
	\end{question}
	\begin{enumerate}
		\item 保温杯结构：保温杯由内外两层薄壁不锈钢筒焊接而成，内外筒之间抽真空以消除传导漏热和对流漏热。内胆外壁涂有防辐射层以降低辐射漏热。
		
		\item 内外壁厚度相同：在对保温杯内胆的质量进行估算时，假设其内外壁厚度相同。这种假设简化了实际可能存在的结构差异。
		
		\item 内胆质量估算：通过测量保温杯的解剖结构和实验时的盛水量，估算与水接触的内壁部分的质量作为 \( m_1 \)。这种估算方法没有考虑内胆材料的具体密度和厚度分布，可能导致实际质量与估算值之间存在差异。
		
		\item 粗略估算：对 \( m_1 \) 的估算是比较粗略的，认为其不会超过保温杯总质量的一半，甚至可能只有总质量的三分之一。这种估算忽略了内胆具体的尺寸、形状和材料特性。
		
		\item 固定 \( m_1 \) 的值：为了减少水量和冰量对 \( m_1 \) 的影响，实验中可能会采取措施使杯内壁整体温度与液体温度相同。这里的假设是杯盖内表面的质量与内胆杯口的质量相同，这可能不完全准确。
		
		\item 误差分析：结构差异：实际的保温杯内胆可能在材质、厚度和制造工艺上有所不同，这可能导致热质量的实际值与估算值有显著差异。估算简化：采用总质量的一定比例进行估算可能没有考虑到实际的物理特性，如不锈钢的具体类型、密度和热容。环境因素：实验中可能存在的其他影响因素，如环境温度变化、保温杯材料的热传导特性等，也可能导致误差。
		
		\item 除此之外，还有：
		忽略环境热量交换:保温杯的设计是为了减缓热量流失，但估算模型可能忽略了环境与保温杯之间的 热量交换。实际上，即使使用了保温杯，仍然存在与环境的一些热量交换，例如通过杯口或盖子的开口。
		
		均匀温度分布:估算模型可能假设液体内部温度是均匀的，即整个液体的温度始终等于测得的温度。 然而，在实际情况中，液体内部可能存在温度梯度，尤其是靠近杯壁的部分。
		
		液体性质的恒定性:模型可能简化为假设液体的性质(如比热容)在整个实验过程中是恒定的。实际 上，液体的性质可能会随着温度的变化而变化，这在实际应用中可能引入一些误差。
		
		热传导的单一途径:模型可能只考虑热量通过液体与保温杯的热传导途径，而忽略了其他途径，如对 流和辐射。在某些情况下，这些额外的热量损失可能对实验结果产生一定的影响。
	\end{enumerate}
		
	\begin{question}
		不用搅拌器对冰的熔化时间延长多久（可通过实验回答），对测量精度有多大的影响？
	\end{question}
	基于热力学原理和实验方法的一般知识，我们可以分析搅拌器的缺失可能对实验的影响：
	
	熔化时间的延长：不使用搅拌器时，保温杯内的水温可能不均匀，导致冰块的熔化速率较慢。搅拌有助于促进热量在水中的均匀分布，从而加快冰的熔化过程。
	
	对测量精度的影响：熔化速度减慢可能导致在达到热平衡之前，系统与环境之间有更多的热量交换，这可能导致实验结果的不准确。另外，水温的不均匀分布也可能导致在测量温度时出现误差。
	
	\begin{question}
		外推法修正漏热的物理原理是什么？能否推导出来？
	\end{question}
	
	外推法修正漏热（extrapolation method for correcting heat leak）是一种用于准确测量热流的方法，特别是在低温物理实验中。这种方法的物理原理基于对实验系统中不可避免的热漏（heat leak）进行量化和修正。热漏是指不受控制的热量交换，通常是由于实验系统与外界环境之间的热传导或辐射造成的。
	
	外推法的基本思想是在不同的已知条件下测量热流，然后通过外推这些测量结果以估计在理想（无漏热）条件下的热流。这通常涉及在不同的温度或功率水平下进行一系列实验，然后分析这些数据以得出漏热的影响。
	
	让我们以一种简化的方式来推导外推法修正漏热的基本方程。假设系统的热流量 \( Q \) 由两部分组成：一部分是由于实验操作（如加热）引起的，另一部分是热漏 \( Q_{\text{leak}} \)。如果我们在不同的操作功率 \( P \) 下测量热流 \( Q \)，我们可以记录 \( Q \) 随 \( P \) 的变化。
	
	我们可以假设热漏 \( Q_{\text{leak}} \) 与外部温度和系统温度之间的差异有关，但在简化模型中，我们可以认为它是恒定的。因此，总热流量 \( Q \) 可以表示为：
	\[ Q = P + Q_{\text{leak}} \]
	
	通过在不同的 \( P \) 下测量 \( Q \)，我们可以绘制 \( Q \) 对 \( P \) 的图表。理想情况下（无热漏），这条线将通过原点。但实际上，由于热漏，这条线会与 \( P \) 轴相交于 \( -Q_{\text{leak}} \)。
	
	\begin{figure}[htbp]
		\centering
		\subfloat[]{
			\includegraphics[width=0.5\textwidth]{AC4GraAD3.png}
		}
		\subfloat[]{
			\includegraphics[width=0.5\textwidth]{AC4GraAD4.png}
		}
		\caption{外推法示意图}
		\label{fig:figAD3}			
	\end{figure}
	
	具体来说，冰的迅速熔解：假设在某一时刻 \( t' \)，冰迅速熔解，系统温度从 \( T'_0 \) 降到 \( T'_1 \)，形成变温曲线中的一个竖直的“熔解线”。漏热的外推：即使系统一直在漏热，通过将单相区（只有水或只有冰的状态）漏热造成的变温曲线外推到与熔解线相交，就可以确定 \( T'_0 \) 和 \( T'_1 \)。这种方法不需要知道环境温度。面积法：具体作图时，熔解线（EF）与变温曲线交于点C。通过移动EF线，使得由线段BCE（蓝色阴影部分）围成的面积等于由线段CFG（红色阴影部分）围成的面积。在这个条件下，初始温度 \( T'_0 \) 对应于熔解线（EF）与降温线（AB）延长线的交点E所对应的温度，结束温度 \( T'_1 \) 对应于熔解线（EF）与升温线（GD）延长线的交点F所对应的温度。
	
	这种方法利用了温度变化的面积法则，即在理想情况下，漏热造成的温度变化所对应的面积应该在整个过程中保持平衡。因此，可以通过调整和比较不同阶段的温度变化曲线的面积，来推导出实际的初始和结束温度。外推法修正漏热的这种方法在理论上是可行的，但它需要精确的温度记录和精细的作图方法。实验中的误差可能来源于温度测量的不准确、曲线作图的不精确、以及实际过程中冰熔化的速率不是瞬时的等因素。
	
	\begin{question}
		如何判断冰的温度为0摄氏度？
	\end{question}
	观察温度计读数，如果温度计示数保持在0摄氏度一段时间，可认为冰的温度为0摄氏度；肉眼观察冰的物态变化，观察到冰开始融化时可认为冰的温度为摄氏度。
	
	%实验记录	
	\clearpage
	\begin{table}
		\renewcommand\arraystretch{1.7}
		\centering
		\begin{tabularx}{\textwidth}{|X|X|X|X|}
			\hline
			专业： & 物理学 & 年级： & 2022级 \\
			\hline
			姓名： & 杨舒云 & 学号： & 22344020\\
			\hline
			室温： & 22摄氏度 & 实验地点： & 实验室515 \\
			\hline
			学生签名：& 杨舒云 & 评分： &\\
			\hline
			实验时间：& 2023/12/7 & 教师签名：&\\
			\hline
		\end{tabularx}
	\end{table}
	
	\section{AC4 冰的熔化热测量 \quad\heiti 实验记录}
	\subsection{实验内容、步骤与结果}
	使用混合法，对已知质量、比热的系统，加入已知质量的冰，测量前后的温度，计算冰的熔化热。
	
	实验会用到的参数：
	
	冰的比热：$C_0=4.18kJ/(kg·^\circ C)$
	
	冰的熔化热（标准值）：$L=334kJ/kg$
	
	\textbf{实验内容}如图所示。
	
	\begin{figure}[htbp]
		\centering
		\includegraphics[width=0.7\textwidth]{AC4GraAD5.jpg}
		\caption{实验内容}
		\label{fig:figAD5}
	\end{figure}
	
	\clearpage
	\subsubsection{用保温杯测量冰的熔化热}
	
	实验数据记录如\cref{tab:tab1}所示。
	
	\begin{table}[h]
		\centering
		\begin{tabular}{|c|c|c|c|c|c|}
			\hline
			测量次序 & $m_c/g$ & $(m_c+m_0)/g$ & $(m_c+m_0+m)/g$ & $T_0/^\circ C$ & $T/^\circ C$ \\
			\hline
			1 & 285.44 & 632.92 & 655.18 & 28.3 & 22.1 \\
			2 & 286.06 & 602.10 & 659.22 & 29.2 & 13.4 \\
			3 & 286.51 & 728.01 & 765.50 & 32.7 & 24.0 \\
			4 & 286.51 & 702.47 & 756.90 & 25.8 & 14.0 \\
			5 & 286.51 & 747.33 & 803.74 & 39.4 & 26.6 \\
			\hline
		\end{tabular}
		\caption{保温杯质量（$m_c$）、水质量（$m_0$）、冰质量（$m$）、初温、末温测量记录}
		\label{tab:tab1}
	\end{table}
	
	对应的由电脑记录的原始数据展示，如\cref{fig:figAA1}所示。
	
	\begin{figure}[htbp]
		\centering
		\subfloat[]{
			\includegraphics[width=0.5\textwidth]{AC4GraAA1.jpg}
		}
		\subfloat[]{
			\includegraphics[width=0.5\textwidth]{AC4GraAA2.jpg}
		}
		\caption{原始数据：用保温杯测量冰的熔化热}
		\label{fig:figAA1}			
	\end{figure}
	
	\clearpage
	\subsubsection{用补偿法（或外推法）在保温杯内实验提高测量精度}
	记录保温杯质量（$m_c$）、水质量（$m_0$）、冰质量（$m$）：（室温 $T =21.5^\circ C$）如\cref{tab:tab2}所示。
	
	\begin{table}[h]
		\centering
		\begin{tabular}{|c|c|c|c|c|}
			\hline
			测量次序 & $m_c/g$ & $(m_c+m_0)/g$ & $(m_c+m_0+m)/g$ & 文件编号 \\
			\hline
			1 & 285.44 & 703.39 & 749.56 & 222 \\
			2 & 285.44 & 688.98 & 752.96 & 21 \\
			3 & 285.44 & 666.38 & 740.33 & 22 \\
			4 & 285.44 & 647.16 & 711.92 & 23 \\
			\hline
		\end{tabular}
		\caption{质量数据记录：用补偿法（或外推法）在保温杯内实验提高测量精度}
		\label{tab:tab2}
	\end{table}
	
	温度-时间序列展示如\cref{tab:tab3}所示（对应文件222）。
	
	\begin{table}[h]
		\centering
		\begin{tabular}{|c|c|c|c|c|c|}
			\hline
			Time/s & Temp/$^\circ C$ & Time/s & Temp/$^\circ C$ & Time/s & Temp/$^\circ C$ \\
			\hline
			0 & 28.6 & 90 & 18.1 & 180 & 17.6 \\
			10 & 26.9 & 110 & 17.9 & 190 & 17.6 \\
			20 & 22.6 & 120 & 17.8 & 200 & 17.7 \\
			40 & 20.9 & 130 & 17.8 & 220 & 17.6 \\
			50 & 20.2 & 140 & 17.7 & 230 & 17.6 \\
			60 & 19.8 & 150 & 17.7 & 240 & 17.7 \\
			70 & 19.2 & 160 & 17.7 & 250 & 17.7 \\
			80 & 18.4 & & & & \\
			\hline
		\end{tabular}
		\caption{时间-温度数据记录：用补偿法（或外推法）在保温杯内实验提高测量精度}
		\label{tab:tab3}
	\end{table}
	
	对应的由电脑记录的原始数据展示，如\cref{fig:figAA3}所示。
	
	\begin{figure}[htbp]
		\centering
		\includegraphics[width=0.5\textwidth]{AC4GraAA3.jpg}
		\caption{原始数据：用补偿法（或外推法）在保温杯内实验提高测量精度}
		\label{fig:figAA3}
	\end{figure}

    \clearpage
	\subsubsection{用补偿法（或外推法）在量热器内实验提高测量精度}
	记录量热器质量（$m_c$）、水质量（$m_0$）、冰质量（$m$）：（室温 $T =21.5^\circ C$）如\cref{tab:tabA4}所示。
	
	\begin{table}[h]
		\centering
		\begin{tabular}{|c|c|c|c|}
			\hline
			测量次序 & $m_c/g$ & $(m_c+m_0)/g$ & $(m_c+m_0+m)/g$ \\
			\hline
			1 & 30.27 & 163.3 & 181.85 \\
			2 & 26.7 & 157.94 & 176.09 \\
			3 & 30.27 & 155.5 & 174.05 \\
			\hline
		\end{tabular}
		\caption{质量数据记录：用补偿法（或外推法）在量热器内实验提高测量精度}
		\label{tab:tabA4}
	\end{table}
	
	温度-时间序列展示如\cref{tab:tab4}所示（对应文件333）。
	
	\begin{table}[h]
		\centering
		\begin{tabular}{|c|c|c|c|c|c|}
			\hline
			Temp/$^\circ C$ & Time/s & Temp/$^\circ C$ & Time/s & Temp/$^\circ C$ & Time/s \\
			\hline
			31.5 & 0 & 19.1 & 60 & 19 & 110 \\
			29 & 10 & 19.1 & 60 & 19 & 120 \\
			28.6 & 10 & 19 & 70 & 19 & 120 \\
			25.7 & 20 & 18.9 & 80 & 19 & 130 \\
			25.5 & 20 & 18.9 & 80 & 19 & 130 \\
			22.6 & 30 & 18.9 & 90 & 19.1 & 140 \\
			22.4 & 30 & 19 & 100 & 19 & 150 \\
			21.2 & 40 & 19 & 100 & 19 & 150 \\
			21.2 & 40 & 19 & 110 & 19.1 & 160 \\
			20.3 & 50 & & & & \\
			\hline
		\end{tabular}
		\caption{时间-温度数据记录：用补偿法（或外推法）在量热器内实验提高测量精度}
		\label{tab:tab4}
	\end{table}	
	
	对应的由电脑记录的原始数据展示，如\cref{fig:figAA4}所示。
	
	\begin{figure}[htbp]
		\centering
		\includegraphics[width=0.4\textwidth]{AC4GraAA4.jpg}
		\caption{原始数据：用补偿法（或外推法）在量热器内实验提高测量精度}
		\label{fig:figAA4}
	\end{figure}
	
	\clearpage
	\subsubsection{数据展示}
	说明：由于主要实验数据由电脑存储为txt文件，因此我采用\textbf{如附录所示的代码}处理数据，以下展示文件“222.txt”进行数据整理后得到的温度-时间序列：
	
	[28.6, 28.6, 28.6, 28.6, 28.6, 28.6, 28.5, 28.5, 28.4, 28.4, 28.3, 28.2, 28.1, 28.0, 27.8, 27.6, 27.3, 26.9, 26.4, 25.9, 25.4, 24.9, 24.5, 24.0, 23.6, 23.2, 22.9, 22.7, 22.6, 22.6, 22.6, 22.6, 22.5, 22.5, 22.6, 22.6, 22.7, 22.7, 22.6, 22.0, 22.1, 22.1, 22.1, 22.1, 22.0, 22.0, 21.9, 21.9, 21.8, 21.7, 21.6, 21.5, 21.4, 21.4, 21.3, 21.2, 21.2, 21.2, 21.1, 21.1, 21.1, 21.1, 21.0, 21.0, 20.9, 20.9, 20.9, 20.8, 20.8, 20.8, 20.8, 20.8, 20.7, 20.7, 20.7, 20.6, 20.6, 20.5, 20.5, 20.4, 20.3, 20.3, 20.2, 20.2, 20.2, 20.1, 20.1, 20.1, 20.1, 20.1, 20.1, 20.1, 20.0, 20.0, 20.0, 20.0, 19.9, 19.9, 19.9, 19.8, 19.8, 19.7, 19.7, 19.6, 19.6, 19.6, 19.5, 19.5, 19.4, 19.4, 19.4, 19.3, 19.3, 19.3, 19.2, 19.2, 19.1, 19.1, 19.1, 19.0, 19.0, 19.0, 18.9, 18.9, 18.9, 18.8, 18.8, 18.7, 18.7, 18.6, 18.5, 18.5, 18.4, 18.4, 18.4, 18.3, 18.3, 18.3, 18.2, 18.2, 18.2, 18.2, 18.2, 18.2, 18.2, 18.2, 18.1, 18.1, 18.1, 18.1, 18.1, 18.1, 18.0, 18.0, 18.0, 18.0, 18.0, 17.9, 17.9, 17.9, 17.9, 17.9, 17.9, 17.9, 17.9, 17.9, 17.9, 17.9, 17.9, 17.9, 17.9, 17.9, 17.9, 17.9, 17.9, 17.9, 17.9, 17.9, 17.9, 17.9, 17.9, 17.9, 17.9, 17.9, 17.9, 17.9, 17.8, 17.8, 17.8, 17.8, 17.8, 17.8, 17.8, 17.8, 17.8, 17.8, 17.8, 17.8, 17.8, 17.8, 17.8, 17.8, 17.8, 17.8, 17.8, 17.8, 17.8, 17.8, 17.8, 17.8, 17.8, 17.8, 17.8, 17.8, 17.8, 17.8, 17.8, 17.8, 17.8, 17.8, 17.8, 17.8, 17.8, 17.8, 17.8, 17.8, 17.7, 17.7, 17.7, 17.7, 17.7, 17.7, 17.7, 17.7, 17.7, 17.7, 17.7, 17.7, 17.7, 17.7, 17.7, 17.7, 17.7, 17.7, 17.7, 17.7, 17.7, 17.7, 17.7, 17.7, 17.7, 17.7, 17.7, 17.7, 17.6, 17.6, 17.7, 17.6, 17.7, 17.7, 17.6, 17.6, 17.6, 17.7, 17.7, 17.6, 17.6, 17.6, 17.6, 17.6, 17.6, 17.6, 17.6, 17.6, 17.6, 17.6, 17.6, 17.6, 17.6, 17.7, 17.7, 17.7, 17.7, 17.7, 17.7, 17.7, 17.7, 17.6, 17.6, 17.6, 17.6, 17.6, 17.6, 17.6, 17.6, 17.6, 17.6, 17.6, 17.6, 17.6, 17.6, 17.6, 17.6, 17.6, 17.6, 17.6, 17.6, 17.6, 17.6, 17.6, 17.6, 17.6, 17.6, 17.6, 17.6, 17.6, 17.6, 17.6, 17.6, 17.6, 17.6, 17.6, 17.6, 17.7, 17.6, 17.6, 17.7, 17.6, 17.6, 17.7, 17.7, 17.7, 17.7, 17.6, 17.6, 17.6, 17.6, 17.6, 17.6, 17.6, 17.6, 17.6, 17.6, 17.6, 17.6, 17.6, 17.6, 17.6, 17.7, 17.7, 17.7, 17.7, 17.7, 17.7, 17.7, 17.7, 17.7, 17.7, 17.7, 17.7, 17.7, 17.6, 17.6, 17.6, 17.6, 17.6, 17.6, 17.6, 17.6, 17.6, 17.6, 17.6, 17.6, 17.6, 17.6, 17.6, 17.6, 17.6, 17.6, 17.6, 17.6, 17.6, 17.6, 17.6, 17.7, 17.7, 17.7, 17.7, 17.7, 17.7, 17.6, 17.7, 17.7, 17.7, 17.7, 17.7, 17.7, 17.7, 17.6, 17.6, 17.6, 17.6, 17.6, 17.6, 17.6, 17.6, 17.6, 17.6, 17.6, 17.6, 17.6, 17.6, 17.7, 17.7, 17.7, 17.7, 17.7, 17.7, 17.7, 17.6, 17.6, 17.6, 17.6, 17.6, 17.6, 17.6, 17.6, 17.6, 17.6, 17.7, 17.7, 17.6, 17.6, 17.6, 17.6, 17.6, 17.6]
	
	[60877, 60877, 60878, 60879, 60879, 60880, 60880, 60881, 60882, 60882, 60883, 60883, 60884, 60885, 60885, 60886, 60886, 60887, 60888, 60888, 60889, 60889, 60890, 60891, 60891, 60892, 60893, 60893, 60894, 60894, 60895, 60896, 60896, 60897, 60897, 60898, 60899, 60899, 60900, 60900, 60901, 60902, 60902, 60903, 60903, 60904, 60905, 60905, 60906, 60906, 60907, 60908, 60908, 60909, 60910, 60910, 60911, 60911, 60912, 60913, 60913, 60914, 60914, 60915, 60916, 60916, 60917, 60917, 60918, 60919, 60919, 60920, 60920, 60921, 60922, 60922, 60923, 60924, 60924, 60925, 60925, 60926, 60927, 60927, 60928, 60928, 60929, 60930, 60930, 60931, 60931, 60932, 60933, 60933, 60934, 60934, 60935, 60936, 60936, 60937, 60937, 60938, 60939, 60939, 60940, 60941, 60941, 60942, 60942, 60943, 60944, 60944, 60945, 60945, 60946, 60947, 60947, 60948, 60948, 60949, 60950, 60950, 60951, 60951, 60952, 60953, 60953, 60954, 60955, 60955, 60956, 60956, 60957, 60958, 60958, 60959, 60959, 60960, 60961, 60961, 60962, 60962, 60963, 60964, 60964, 60965, 60965, 60966, 60967, 60967, 60968, 60969, 60969, 60970, 60970, 60971, 60972, 60972, 60973, 60973, 60974, 60975, 60975, 60976, 60976, 60977, 60978, 60978, 60979, 60979, 60980, 60981, 60981, 60982, 60982, 60983, 60984, 60984, 60985, 60986, 60986, 60987, 60987, 60988, 60989, 60989, 60990, 60990, 60991, 60992, 60992, 60993, 60993, 60994, 60995, 60995, 60996, 60996, 60997, 60998, 60998, 60999, 61000, 61000, 61001, 61001, 61002, 61003, 61003, 61004, 61004, 61005, 61006, 61006, 61007, 61007, 61008, 61009, 61009, 61010, 61010, 61011, 61012, 61012, 61013, 61013, 61014, 61015, 61015, 61016, 61017, 61017, 61018, 61018, 61019, 61020, 61020, 61021, 61021, 61022, 61023, 61023, 61024, 61024, 61025, 61026, 61026, 61027, 61027, 61028, 61029, 61029, 61030, 61031, 61031, 61032, 61032, 61033, 61034, 61034, 61035, 61035, 61036, 61037, 61037, 61038, 61038, 61039, 61040, 61040, 61041, 61041, 61042, 61043, 61043, 61044, 61045, 61045, 61046, 61046, 61047, 61048, 61048, 61049, 61049, 61050, 61051, 61051, 61052, 61052, 61053, 61054, 61054, 61055, 61055, 61056, 61057, 61057, 61058, 61059, 61059, 61060, 61060, 61061, 61062, 61062, 61063, 61063, 61064, 61065, 61065, 61066, 61066, 61067, 61068, 61068, 61069, 61069, 61070, 61071, 61071, 61072, 61073, 61073, 61074, 61074, 61075, 61076, 61076, 61077, 61077, 61078, 61079, 61079, 61080, 61080, 61081, 61082, 61082, 61083, 61083, 61084, 61085, 61085, 61086, 61086, 61087, 61088, 61088, 61089, 61090, 61090, 61091, 61091, 61092, 61093, 61093, 61094, 61094, 61095, 61096, 61096, 61097, 61097, 61098, 61099, 61099, 61100, 61100, 61101, 61102, 61102, 61103, 61104, 61104, 61105, 61105, 61106, 61107, 61107, 61108, 61108, 61109, 61110, 61110, 61111, 61111, 61112, 61113, 61113, 61114, 61114, 61115, 61116, 61116, 61117, 61118, 61118, 61119, 61119, 61120, 61121, 61121, 61122, 61122, 61123, 61124, 61124, 61125, 61125, 61126, 61127, 61127, 61128, 61128, 61129, 61130, 61130, 61131, 61131, 61132, 61133, 61133, 61134, 61135, 61135, 61136, 61136, 61137, 61138, 61138, 61139, 61139, 61140, 61141, 61141, 61142]
	
	\begin{figure}[htbp]
		\centering
		\includegraphics[width=0.7\textwidth]{AC4Gra1.jpg}
		\caption{数据处理界面}
		\label{fig:fig1}
	\end{figure}
	
	\clearpage
	\subsection{原始数据记录}
	实验记录本上的原始数据见\cref{fig:figA1}以及由电脑存储的原始数据见\cref{fig:figA2}。
	
	\begin{figure}[htbp]
		\centering
		\includegraphics[width=0.7\textwidth]{AC4GraA1.png}
		\caption{原始数据记录}
		\label{fig:figA1}
	\end{figure}
	
	\begin{figure}[htbp]
		\centering
		\includegraphics[width=0.58\textwidth]{AC4GraA2.jpg}
		\caption{原始数据记录}
		\label{fig:figA2}
	\end{figure}
	
	\clearpage
	实验过程中的操作界面见\cref{fig:figA3}。
	
	\begin{figure}[htbp]
		\centering
		\subfloat[]{
			\includegraphics[width=0.3\textwidth]{AC4GraA3.jpg}
		}
		\subfloat[]{
			\includegraphics[width=0.3\textwidth]{AC4GraA4.jpg}
		}
		\caption{实验过程中的操作界面}
		\label{fig:figA3}			
	\end{figure}
	
	实验台桌面整理见\textbf{附件}部分(\cref{fig:figA5})。
	
	\subsection{实验过程中遇到的问题记录}
	\begin{enumerate}
		\item 测不准！
		\item 冰块无法一次性放入。
		\item 热电偶碰到冰块和杯壁都会造成较大误差
	\end{enumerate}
	
	
	
	%分析与讨论	
	\clearpage
	\begin{table}
		\renewcommand\arraystretch{1.7}
		\begin{tabularx}{\textwidth}{|X|X|X|X|}
			\hline
			专业：& 物理学 &年级：& 2022级\\
			\hline
			姓名： & 杨舒云 & 学号：& 22344020\\
			\hline
			日期：& 2023/12/7 & 评分： &\\
			\hline
		\end{tabularx}
	\end{table}
	
	\section{AC4 冰的熔化热测量 \quad\heiti 分析与讨论}
	\subsection{实验数据分析}
	
	\subsubsection{计算}
	水、冰质量计算：
	\[m_0=m_c+m_0-m_c\]
	\[m=m_c+m_0+m-m_c-m_0\]
	
	潜热计算：
	\[L=1/m (m_0 C_0+m_1 C_1 )(T_0-T_1 )-C_0(T_1-T')\]
	
	实验1（用保温杯测量冰的熔化热）计算得到如\cref{tab:tab5}所示结果。
	
	\begin{table}[h]
		\centering
		\begin{tabular}{|c|c|c|c|c|c|c|c|}
			\hline
			$m_c/g$ & $(m_c+m_0)/g$ & $(m_c+m_0+m)/g$ & $T_0/^\circ C$ & $T/^\circ C$ & $m/g$ & $m_0/g$ & $L/(J/g)$ \\
			\hline
			285.44 & 632.92 & 655.18 & 28.30 & 22.10 & 22.26 & 347.48 & 324.35 \\
			286.06 & 602.10 & 659.22 & 29.20 & 13.40 & 57.12 & 316.04 & 321.52 \\
			286.51 & 728.01 & 765.50 & 32.70 & 24.00 & 37.49 & 441.50 & 338.13 \\
			286.51 & 702.47 & 756.90 & 25.80 & 14.00 & 54.43 & 415.96 & 327.93 \\
			286.51 & 747.33 & 803.74 & 39.40 & 26.60 & 56.41 & 460.82 & 335.85 \\
			\hline
		\end{tabular}
		\caption{潜热计算}
		\label{tab:tab5}
	\end{table}	
	
	\clearpage
	\subsubsection{（仅对实验内容2）根据第5步至第7步所记录的数据，在二维直角坐标上作出对应的温度-时间曲线，并作出B点(投冰)、G点（温度回升点）对应的温度值分别为$T_0$和$T_1$}
	
	在实验2（用补偿法（或外推法）在保温杯内实验提高测量精度）中我和同学戴鹏辉22344016发现了一种小妙招，即\textbf{用稿纸折成漏斗辅助投冰}，可使得冰尽可能于同一时间投入保温杯中。
	
	利用实验所得数据，作图如下所示（部分数据利用附件中得代码直接作图，部分数据利用Origin作更美观的图）。
	\begin{enumerate}
		\item 对应文件“222.txt”如\cref{fig:fig2}所示。
		
		B：$28.6^\circ C$；G：$17.6^\circ C$。
		
		\begin{figure}[htbp]
			\centering
			\includegraphics[width=0.6\textwidth]{AC4Gra2.jpg}
			\caption{温度-时间曲线}
			\label{fig:fig2}
		\end{figure}
		
		\item 对应文件“22.txt”如\cref{fig:fig3}所示。
		
		B：$31.5^\circ C$；G：$13.4^\circ C$。
		
		\begin{figure}[htbp]
			\centering
			\includegraphics[width=0.6\textwidth]{AC4Gra3.jpg}
			\caption{温度-时间曲线}
			\label{fig:fig3}
		\end{figure}
		
		\item \cref{fig:fig4}来自于“23.txt”和“333.txt”（实验3：用补偿法（或外推法）在量热器内实验提高测量精度）。
		
		B：$30.5^\circ C$；G：$14.5^\circ C$。
		
		B：$31.5^\circ C$；G：$18.9^\circ C$。
		
		\begin{figure}[htbp]
			\centering
			\subfloat[]{
				\includegraphics[width=0.5\textwidth]{AC4Gra4.png}
			}
			\subfloat[]{
				\includegraphics[width=0.5\textwidth]{AC4Gra5.png}
			}
			\caption{温度-时间曲线}
			\label{fig:fig4}			
		\end{figure}
		
	\end{enumerate}
	
	\clearpage
	\subsubsection{求出冰的熔化热，与标准值334J/g相比较}
	\begin{enumerate}
		\item 实验1计算得到熔化热及相对误差如\cref{tab:tab6}所示。
		
		\begin{table}[h]
			\centering
			\begin{tabular}{|c|c|c|c|c|c|c|c|c|}
				\hline
				$m_c/g$ & $(m_c+m_0)/g$ & $(m_c+m_0+m)/g$ & $T_0/^\circ C$ & $T/^\circ C$ & $m/g$ & $m_0/g$ & $L/(J/g)$ & 相对误差 \\
				\hline
				285.44 & 632.92 & 655.18 & 28.30 & 22.10 & 22.26 & 347.48 & 324.35 & -2.89\% \\
				286.06 & 602.10 & 659.22 & 29.20 & 13.40 & 57.12 & 316.04 & 321.52 & -3.74\% \\
				286.51 & 728.01 & 765.50 & 32.70 & 24.00 & 37.49 & 441.50 & 338.13 & 1.24\% \\
				\hline
			\end{tabular}
			\caption{实验1计算得到熔化热及相对误差}
			\label{tab:tab6}
		\end{table}		
		
		\item 实验2利用\textbf{附录中的代码}采用补偿法及外推法计算得到熔化热及相对误差如\cref{tab:tab7}所示。
		
		\begin{table}[h]
			\centering
			\begin{tabular}{|c|c|c|c|c|c|c|c|c|}
				\hline
				$m_c/g$ & $(m_c+m_0)/g$ & $(m_c+m_0+m)/g$ & $T_0/^\circ C$ & $T/^\circ C$ & $m/g$ & $m_0/g$ & $L/(J/g)$ & 相对误差 \\
				\hline
				285.44 & 703.39 & 749.56 & 28.6 & 17.9 & 46.17 & 417.95 & 340.19 & 1.85\% \\
				285.44 & 666.38 & 740.33 & 31.5 & 15.5 & 73.95 & 380.94 & 308.14 & -7.74\% \\
				285.44 & 647.16 & 711.92 & 30.50 & 17.1 & 64.76 & 361.72 & 268.55 & -19.60\% \\
				\hline
			\end{tabular}
			\caption{实验2利用\textbf{附录中的代码}采用补偿法及外推法计算得到熔化热及相对误差}
			\label{tab:tab7}
		\end{table}
		
		\item 实验3利用\textbf{附录中的代码}采用补偿法及外推法计算得到熔化热及相对误差如\cref{tab:tab8}所示。
		
		\begin{table}[h]
			\centering
			\begin{tabular}{|c|c|c|c|c|c|c|c|c|}
				\hline
				$m_c/g$ & $(m_c+m_0)/g$ & $(m_c+m_0+m)/g$ & $T_0/^\circ C$ & $T/^\circ C$ & $m/g$ & $m_0/g$ & $L/(J/g)$ & 相对误差 \\
				\hline
				30.27 & 163.30 & 181.85 & 31.50 & 19 & 18.55 & 133.03 & 304.67 & -8.78\% \\
				27.73 & 158.44 & 178.66 & 33.70 & 20.10 & 20.22 & 130.71 & 292.05 & -12.56\% \\
				\hline
			\end{tabular}
			\caption{实验3利用\textbf{附录中的代码}采用补偿法及外推法计算得到熔化热及相对误差}
			\label{tab:tab8}
		\end{table}
		
	\end{enumerate}
	
	\subsubsection{误差分析}
	
	\begin{enumerate}
		\item 我们已经计算了相对误差，现将结果陈列于\cref{tab:tab9}。其中，1、2和3来自于实验1，4、5和6来自于实验2，7和8来自于实验3.
		
		\begin{table}[h]
			\centering
			\begin{tabular}{|c|c|c|}
				\hline
				测量序号 & L/(J/g) & 相对误差 \\
				\hline
				1 & 324.35 & -2.89\% \\
				2 & 321.52 & -3.74\% \\
				3 & 338.13 & 1.24\% \\
				4 & 340.19 & 1.85\% \\
				5 & 308.14 & -7.74\% \\
				6 & 268.55 & -19.60\% \\
				7 & 304.67 & -8.78\% \\
				8 & 292.05 & -12.56\% \\
				\hline
			\end{tabular}
			\caption{相对误差}
			\label{tab:tab9}
		\end{table}
		
		从相对误差表可以看出，总得来说，实验1的误差是最小的。
		
		\item 误差来源
		
		温度测量误差：温度是混合法测量冰的熔化热的重要参数之一。由于温度计的精度和准确性，环境温度的变化，以及温度读数的误差，都会影响温度测量值的准确性。热电偶温度计可能会碰到冰块，使得示数不准。
		
		质量测量误差：混合法测量冰的熔化热需要用到已知质量的物质，如水，内筒，搅拌器等。如果称量这些物质的质量不准确，或者在实验过程中有水溅出或蒸发，都会影响质量测量值的准确性。
		
		混合过程误差：操作过程中冰块很难一次性放入杯中， 则增加了冰块与环境交换热量的时间，导致误差增大。混合过程中需要保证混合物质充分混合，且冰完全熔化。如果混合不均匀，或者冰未完全熔化，或者冰的纯度不高，都会影响混合过程的准确性。
		
		散热修正误差：测量时间不应过长，会增加与外界交换的热量，增大误差。由于量热器的绝热条件并不完善，实验系统并非严格的孤立系统，所以在实验过程中，系统会与外界交换热量，导致测量结果偏小。为了修正这种散热误差，可以采用补偿法，即使系统的初温和末温分别高于和低于环境温度，使系统与外界交换的热量前后相抵消。
		
		综上所述，混合法测量冰的熔化热的误差主要来源于温度测量误差，质量测量误差，混合过程误差，和散热修正误差等方面。为了减小误差，提高测量结果的准确性，需要注意以下几点：选择合适的温度计，定期校准，准确读数，记录环境温度；选择合适的电子天平，准确称量，避免水溅出或蒸发；选择合适的冰块，擦干表面，快速投入，充分搅拌，确保完全熔化；选择合适的水温，使初温和末温分别高于和低于环境温度，达到散热补偿的效果。
		
		\item 系统误差和随机误差
		
		该实验的系统误差主要包括以下几个方面：
		
		温度计的系统误差：温度计本身的刻度或校准可能不准确，导致测量温度与真实温度有一定的偏差。这种误差可以通过对温度计进行校准或换用更精确的温度计来避免或减小。观察数据图像可以得知，在降温过程中，存在温度数据波动，这是主要因为热电偶温度计的探头测量的可能并不是完全均匀的温度，也就是说当探头的位置发生改变，就可能导致温度数据发生突变。
		
		冰的系统误差：冰的纯度、温度、表面积等因素可能影响冰的熔化热的测量。这种误差可以通过选择纯度高、温度低、表面积小的冰块，以及擦干冰块表面的水分来避免或减小。
		
		散热的系统误差：由于量热器的绝热条件并不完善，实验系统并非严格的孤立系统，所以在实验过程中，系统会与外界交换热量，导致测量结果偏小。由于实验过程中存在长时间使冰和水直接暴露在空气中的操作（操作问题），导致了大量漏热。实验装置的盖子可能会导致漏热情况的发生（不完全密封）。这种误差可以通过采用补偿法，即使系统的初温和末温分别高于和低于环境温度，使系统与外界交换的热量前后相抵消来避免或减小。
		
		该实验的随机误差主要包括以下几个方面：
		
		温度读数的随机误差：由于温度计的读数精度有限，或者操作者的读数技巧不同，导致测量温度与读数温度有一定的偏差。这种误差可以通过增加测量次数，取平均值来减小。
		
		质量称量的随机误差：由于电子天平的称量精度有限，或者操作者的称量技巧不同，导致测量质量与称量质量有一定的偏差。这种误差可以通过增加测量次数，取平均值来减小。
		
		混合过程的随机误差：由于混合物质的搅拌速度、时间、均匀度等因素的不确定性，导致混合过程的热量交换不完全或不一致，影响测量结果的准确性。这种误差可以通过控制搅拌的条件，使之尽可能一致来减小。
	\end{enumerate}
	
	\subsubsection{提出实验改进办法}
	\begin{enumerate}
		\item  使用纯度高、温度低、表面积小的冰块：冰的纯度、温度、表面积等因素可能影响冰的熔化热的测量。
		如果冰含有杂质，会降低冰的熔点，使冰的熔化热偏小。如果冰的温度高于0℃，会使冰的熔化热偏大。如果冰的表面积大，会增加冰与水的接触面积，加快冰的熔化速度，使系统温度变化不稳定。因此，应该选择纯度高、温度低、表面积小的冰块，以减小这些因素的影响。
		\item 擦干冰块表面的水分：冰块表面的水分会影响冰的质量测量和熔化过程。如果冰块表面有水分，会使冰的质量偏大，导致冰的熔化热偏小。如果冰块表面有水分，会使冰的熔化速度加快，使系统温度变化不稳定。因此，应该在投入冰块之前，用干净的纸巾擦干冰块表面的水分，以减小这些因素的影响。
		\item 快速投入冰块，充分搅拌，确保完全熔化：冰块的投入速度、搅拌程度、熔化程度等因素也会影响冰的熔化热的测量。如果冰块投入速度慢，会使冰在空气中散热，导致冰的熔化热偏大。如果搅拌不充分，会使混合物质的温度不均匀，导致温度测量误差。如果冰未完全熔化，会使系统温度不稳定，导致热量交换不完全。因此，应该快速投入冰块，充分搅拌，确保完全熔化，以减小这些因素的影响。
		\item 使用外推法测定水的初温：外推法是一种利用系统温度随时间变化的规律，来推算出水的初始温度的方法，它可以避免由于水倒入量热器时的热损失而导致的温度测量误差。具体操作方法是，在倒入水之前，先测量量热器内的空气温度，然后倒入水，立即开始记录系统温度随时间的变化，绘制温度-时间曲线，用外推法求出水的初始温度。
		\item 使用散热补偿法修正散热误差：散热补偿法是一种利用牛顿冷却定律，使系统与外界交换的热量前后相抵消，从而修正散热误差的方法，它可以避免由于量热器的绝热条件不完善而导致的测量结果偏小的问题。具体操作方法是，选择合适的水温，使系统的初温和末温分别高于和低于环境温度，测量系统温度随时间的变化，绘制温度-时间曲线，用图解法求出系统与外界交换的热量，进行散热修正。
	\end{enumerate}	
	
	\clearpage
	\subsection{实验后思考题}
	\begin{question}
		往量热器内投入冰块后，冰块浮在水面，量热器内温度不均匀，对实验有何影响？（针对所采用实验方案进行讨论）。
	\end{question}
	\begin{itemize}
		\item 影响系统的热平衡条件：如果冰块浮在水面，那么冰块与水的接触面积减小，导致冰的熔化速度减慢，使系统达到热平衡的时间延长。这可能导致实验时间过长，增加了热量损失的可能性。同时，如果量热器内温度不均匀，那么系统的平均温度可能不等于温度计测量的温度，导致温度测量误差。这些因素都会影响系统的热平衡条件，从而影响冰的熔化热的测量结果。
		
		\item 影响外推法和散热补偿法的有效性：如果冰块浮在水面，那么冰块与空气的接触面积增大，导致冰的熔化过程受到空气温度的影响，使系统温度随时间的变化不符合牛顿冷却定律。这可能导致外推法和散热补偿法的计算结果不准确，从而影响冰的熔化热的测量结果。
		
		\item 将冰块浸入水中，而不是让其浮在水面上：这样可以增加冰块与水的接触面积，加快冰的熔化速度，缩短实验时间，减小热量损失的可能性。同时，也可以减小冰块与空气的接触面积，减小空气温度对冰的熔化过程的影响，使系统温度随时间的变化更符合牛顿冷却定律，提高外推法和散热补偿法的计算准确性。
		
		\item 使用搅拌器确保量热器内的温度更加均匀：这样可以消除温度梯度，使系统的平均温度等于温度计测量的温度，减小温度测量误差。
		同时，也可以使系统更快地达到热平衡，缩短实验时间，减小热量损失的可能性。
		
	\end{itemize}
	
	\begin{question}
		除搅拌外，如何进一步减少测量影响？
	\end{question}
	在测量冰的熔化热的实验中，如果冰块浮在水面并且量热器内温度不均匀，会对实验结果产生一些影响。
	如果量热器内温度不均匀，那么温度测量可能不够准确。温度梯度可能导致热量的非均匀分布，从而影响对冰熔化热的准确测量。一旦温度计探头位置改变，也会导致温度发生突变。冰块浮在水面上时，与环境之间的接触面积增大，从而增加了热量的损失。这可能导致实际释放到水中的热量小于理论值，从而影响测量的准确性。同时，冰块与空气之间也可能发生热交换，进一步增加热量损失。冰块浮在水面上时，与水之间的热交换可能不够均匀。这可能导致水的温度变化不均匀，进而影响测量的结果。
	
	措施如下：
	\begin{enumerate}
		\item 将冰块浸入水中，而不是让其浮在水面上：这样可以增加冰块与水的接触面积，加快冰的熔化速度，缩短实验时间，减小热量损失的可能性。同时，也可以减小冰块与空气的接触面积，减小空气温度对冰的熔化过程的影响，使系统温度随时间的变化更符合牛顿冷却定律，提高外推法和散热补偿法的计算准确性。
		\item 使用搅拌器确保量热器内的温度更加均匀：这样可以消除温度梯度，使系统的平均温度等于温度计测量的温度，减小温度测量误差。同时，也可以使系统更快地达到热平衡，缩短实验时间，减小热量损失的可能性。
		\item 使用纯度高、温度低、表面积小的冰块：冰的纯度、温度、表面积等因素可能影响冰的熔化热的测量。如果冰含有杂质，会降低冰的熔点，使冰的熔化热偏小 。如果冰的温度高于0摄氏度，会使冰的熔化热偏大 。如果冰的表面积大，会增加冰与水的接触面积，加快冰的熔化速度，使系统温度变化不稳定 。因此，应该选择纯度高、温度低、表面积小的冰块，以减小这些因素的影响。
		\item 擦干冰块表面的水分：冰块表面的水分会影响冰的质量测量和熔化过程。如果冰块表面有水分，会使冰的质量偏大，导致冰的熔化热偏小 。如果冰块表面有水分，会使冰的熔化速度加快，使系统温度变化不稳定 。因此，应该在投入冰块之前，用干净的纸巾擦干冰块表面的水分，以减小这些因素的影响。
		\item 快速投入冰块，充分搅拌，确保完全熔化：冰块的投入速度、搅拌程度、熔化程度等因素也会影响冰的熔化热的测量。如果冰块投入速度慢，会使冰在空气中散热，导致冰的熔化热偏大 。如果搅拌不充分，会使混合物质的温度不均匀，导致温度测量误差 。如果冰未完全熔化，会使系统温度不稳定，导致热量交换不完全 。因此，应该快速投入冰块，充分搅拌，确保完全熔化，以减小这些因素的影响。
		\item 使用外推法测定水的初温：外推法是一种利用系统温度随时间变化的规律，来推算出水的初始温度的方法，它可以避免由于水倒入量热器时的热损失而导致的温度测量误差。具体操作方法是，在倒入水之前，先测量量热器内的空气温度，然后倒入水，立即开始记录系统温度随时间的变化，绘制温度-时间曲线，用外推法求出水的初始温度。
		\item 使用散热补偿法修正散热误差：散热补偿法是一种利用牛顿冷却定律，使系统与外界交换的热量前后相抵消，从而修正散热误差的方法，它可以避免由于量热器的绝热条件不完善而导致的测量结果偏小的问题。具体操作方法是，选择合适的水温，使系统的初温和末温分别高于和低于环境温度，测量系统温度随时间的变化，绘制温度-时间曲线，用图解法求出系统与外界交换的热量，进行散热修正。
	\end{enumerate}
	
	
	
	\clearpage
	\section{AC4 冰的熔化热测量 \quad\heiti 参考文献与附件}
	\subsection{参考文献}
	[1] 沈韩.基础物理实验.——北京：科学出版社，2015.2 ISBN：978-7-03-043311-4 
	
	[2] 维基百科
	
	[3] 实验讲义
	
	\subsection{附件}
	实验台桌面整理：
	
	\begin{figure}[htbp]
		\centering
		\includegraphics[width=0.7\textwidth]{AC4GraA5.jpg}
		\caption{实验台桌面整理}
		\label{fig:figA5}
	\end{figure}
	
	代码展示：
	
	\begin{lstlisting}[style=pythonstyle,caption=数据处理]
		# 读取txt文件中的数据，假设文件名为data.txt
		with open("333.txt", "r") as f:
		lines = f.readlines()
		
		# 创建一个空列表来存储提取的数据
		data = []
		
		# 遍历每一行数据
		for line in lines:
		# 用空格分隔每一列数据
		columns = line.split()
		#    print(columns[5][0])
		# 提取第三列数据中“/”前的部分，即31.5
		temp = columns[3].split("/")[0]
		# 提取第五列数据中的时分秒部分，即17:18:40
		time = columns[5].split("\\")[1]
		# 将时分秒格式化为以秒为单位的格式，即17*3600 + 18*60 + 40 = 62200
		seconds = sum(int(x) * 60 ** i for i, x in enumerate(reversed(time.split(":"))))
		# 将提取的数据添加到列表中，作为一个元组
		data.append((temp, seconds))
		
		# 创建一个空列表来存储每十秒抽取的样本
		samples = []
		
		# 设置一个初始时间，假设为第一行数据的时间
		initial_time = data[0][1]
		
		# 遍历提取的数据
		for d in data:
		# 如果数据的时间与初始时间的差值是10的倍数，说明是每十秒的样本
		if (d[1] - initial_time) % 10 == 0:
		# 将样本添加到列表中
		samples.append(d)
		
		# 打印每十秒抽取的样本
		#print(samples)
		#print(data)
		
		temp_list = []
		time_list = []
		
		# 遍历每个元组
		for t in data:
		# 将第一个数据添加到温度列表中
		#    temp_list.append(float(t[0]) - 21.2)
		temp_list.append(float(t[0]) - 21.4)
		# 将第二个数据添加到时间列表中
		time_list.append(t[1])
		
		# 打印新的两列数据
		print(temp_list)
		print(time_list)
		
		tempS = []
		timeS = []
		
		for tt in samples:
		tempS.append(float(tt[0]))
		timeS.append(tt[1])
		
		#print(tempS)
		#print(timeS)
		
		# 导入pandas库，用于操作Excel文件
		import pandas as pd
		
		# 创建一个空的DataFrame对象，用于存储数据
		df = pd.DataFrame()
		
		# 将temp_list和time_list作为两列添加到DataFrame中，分别命名为Temp和Time
		df['Temp'] = temp_list
		df['Time'] = time_list
		
		# 创建一个ExcelWriter对象，用于写入Excel文件，假设文件名为data.xlsx
		writer = pd.ExcelWriter('AC4DataOpr.xlsx')
		
		# 将DataFrame中的数据写入Excel文件的第一个工作表，命名为Sheet1
		df.to_excel(writer, 'Sheet1', index=False)
		
		# 保存并关闭Excel文件
		writer.save()
		writer.close()
		
		dfS = pd.DataFrame()
		dfS['Temp'] = tempS
		dfS['Time'] = timeS
		writer = pd.ExcelWriter('AC4DataOpr.xlsx')
		dfS.to_excel(writer, 'Sheet2', index=False)
		writer.save()
		writer.close()
		
		# 导入matplotlib.pyplot库，用于绘制图形
		import matplotlib.pyplot as plt
		
		# 将temp_list中的数据类型转换为浮点数
		temp_list = [float(t) for t in temp_list]
		
		# 绘制折线图，将time_list作为x轴，temp_list作为y轴
		plt.plot(time_list, temp_list)
		
		# 设置x轴和y轴的标签
		plt.xlabel('Time')
		plt.ylabel('Temp')
		
		# 设置图形的标题
		plt.title('Temp vs Time')
		
		# 显示图形
		plt.show()
	\end{lstlisting}
	
	\begin{lstlisting}[style=pythonstyle,caption=补偿法]
		# 导入numpy库，用于数组的处理
		import numpy as np
		from scipy.optimize import bisect
		
		# 导入scipy库，用于数值积分
		from scipy import integrate
		
		# 将温度和时间分别存储为两个numpy数组
		temp_T = np.array(temp_list)
		time_t = np.array(time_list)
		
		# 定义温度-时间曲线函数，使用线性插值法
		def f(t):
		return np.interp(t, time_t, temp_T)
		
		# 定义t0时刻，根据你的实际需求进行修改
		t0 = 62334
		
		# 定义一个函数，计算从t0到t的积分值
		def F(t):
		return integrate.quad(f, t0, t)[0]
		
		# 定义一个函数，计算积分值与0的差值
		def diff(t):
		return F(t) - 0
		
		# 创建一个空列表，用于存储所有使得积分值为0的t值
		t_list = []
		
		# 定义t的取值范围，根据你的实际数据进行修改
		t_min = 62334
		t_max = 62496
		
		# 定义一个循环，从t_min开始，每次增加一个小的步长，直到达到t_max
		# 假设步长为10，你可以根据你的实际数据进行修改
		step = 10
		t = t_min
		while t <= t_max:
		# 使用二分法，找到使得积分值为0的t值，如果存在的话
		try:
		t_zero = bisect(diff, t, t + step)
		# 将找到的t值添加到列表中
		t_list.append(t_zero)
		except ValueError:
		# 如果没有找到t值，忽略错误，继续循环
		pass
		# 更新t值，增加一个步长
		t += step
		
		# 打印结果
		print("The values of t that make the integral zero are:")
		for t in t_list:
		print("{:.2f}".format(t))
		print("The temperature value corresponding to t=62350 is {:.2f}".format(f(62455.51) + 21.4))
	\end{lstlisting}
	
	\begin{lstlisting}[style=pythonstyle,caption=外推法]
		import numpy as np
		import matplotlib.pyplot as plt
		from scipy.optimize import curve_fit
		
		# 示例数据
		time = time_list  # 时间
		temperature = temp_list  # 温度
		
		# 确定冰熔化的时间段
		# 假设熔化发生在温度首次下降到0度的时间点
		melting_start_index = temperature.index(0)
		melting_end_index = len(temperature) - 1 - temperature[::-1].index(0)
		
		# 使用冰熔化前的数据进行线性拟合
		time_before_melting = time[:melting_start_index]
		temp_before_melting = temperature[:melting_start_index]
		
		# 定义线性函数
		def linear_func(x, a, b):
		return a * x + b
		
		# 进行线性拟合
		params, _ = curve_fit(linear_func, time_before_melting, temp_before_melting)
		
		# 使用拟合参数预测冰熔化时的温度
		predicted_temp_at_melting = linear_func(time[melting_start_index], *params)
		predicted_temp_after_melting = linear_func(time[melting_end_index], *params)
		
		# 绘制图形
		plt.figure(figsize=(10, 6))
		plt.plot(time, temperature, label='实际温度')
		plt.plot(time_before_melting + [time[melting_end_index]], 
		[linear_func(t, *params) for t in time_before_melting] + [predicted_temp_after_melting],
		label='外推温度', linestyle='--')
		plt.scatter([time[melting_start_index], time[melting_end_index]], 
		[predicted_temp_at_melting, predicted_temp_after_melting], color='red')
		plt.xlabel('时间')
		plt.ylabel('温度')
		plt.legend()
		plt.title('冰熔化热实验的温度外推法')
		plt.grid(True)
		plt.show()
		
		(predicted_temp_at_melting, predicted_temp_after_melting)
	\end{lstlisting}

\end{document}